\documentclass[11pt, a4paper]{article}

% Gemeinsame Style-Definitionen (TEXINPUTS muss templates/ enthalten — siehe scripts/build_tex.py)
% bfw-style.tex — gemeinsame Style-Definitionen für alle Handouts/Aufgabenhefte
% Voraussetzung: Kompilation mit xelatex (wegen fontspec)

\usepackage[ngerman]{babel}
% Babel-ngerman macht " zu einem aktiven Shorthand (z.B. für "a -> ä). In unseren
% Texten verwenden wir " als normales Zeichen — daher Shorthand komplett ausschalten.
\shorthandoff{"}
\usepackage{geometry}
\geometry{a4paper, top=2.4cm, bottom=2.4cm, left=2.3cm, right=2.3cm,
          headheight=15pt, headsep=0.7cm, footskip=1.1cm}

\usepackage{fontspec}
% Fallback-Kette: Source Sans Pro -> Calibri -> Segoe UI -> Latin Modern Sans
% (Latin Modern ist immer mit MiKTeX vorhanden)
\IfFontExistsTF{Source Sans Pro}{%
  \setmainfont{Source Sans Pro}[Ligatures=TeX]%
}{\IfFontExistsTF{Calibri}{%
  \setmainfont{Calibri}[Ligatures=TeX]%
}{\IfFontExistsTF{Segoe UI}{%
  \setmainfont{Segoe UI}[Ligatures=TeX]%
}{%
  \setmainfont{Latin Modern Sans}[Ligatures=TeX]%
}}}
\IfFontExistsTF{Source Code Pro}{%
  \setmonofont{Source Code Pro}[Scale=0.92]%
}{\IfFontExistsTF{Consolas}{%
  \setmonofont{Consolas}[Scale=0.92]%
}{%
  \setmonofont{Latin Modern Mono}[Scale=0.92]%
}}

% Gesperrte Versalien für Labels/Titelbalken (Kapitälchen-Ersatz, funktioniert
% mit jeder OpenType-Schrift — Latin Modern Sans hat keine echten Kapitälchen)
\newcommand{\bfwlabelfont}{\footnotesize\bfseries\addfontfeatures{LetterSpace=8.0}}

\usepackage{xcolor}
% Farbpalette: hell, freundlich, modern
\definecolor{bfwPrimary}{HTML}{0EA5A4}    % Petrol/Türkis - Primär
\definecolor{bfwPrimaryDark}{HTML}{0F766E} % dunkler Petrol für Headings
\definecolor{bfwAccent}{HTML}{F59E0B}     % Amber - Akzent
\definecolor{bfwSuccess}{HTML}{10B981}    % Grün - Erfolg / Lösung
\definecolor{bfwWarn}{HTML}{F97316}       % Orange - Achtung
\definecolor{bfwInk}{HTML}{0F172A}        % fast Schwarz
\definecolor{bfwInkSoft}{HTML}{475569}    % gedämpftes Grau
\definecolor{bfwBgSoft}{HTML}{F8FAFC}     % off-white
\definecolor{bfwBgTeal}{HTML}{ECFEFF}     % helles Türkis
\definecolor{bfwBgAmber}{HTML}{FEF3C7}    % helles Gelb
\definecolor{bfwBgRose}{HTML}{FFE4E6}     % helles Rosé für Eselsbrücken
\definecolor{bfwTabHead}{HTML}{E4F3F2}    % zarte Kopfzeilen-Hinterlegung für Tabellen

\usepackage[colorlinks=true, linkcolor=bfwPrimaryDark, urlcolor=bfwPrimaryDark]{hyperref}
\usepackage{lastpage}
\usepackage{enumitem}
\setlist[itemize]{leftmargin=*, itemsep=2pt, topsep=2pt}
\setlist[enumerate]{leftmargin=*, itemsep=2pt, topsep=2pt}

\usepackage[most]{tcolorbox}
\usepackage{booktabs}
\usepackage{colortbl}
\usepackage{tikz}
\usetikzlibrary{decorations.pathmorphing, positioning, shapes.geometric, arrows.meta}

% ---------------------------------------------------------------------------
% Einheitliches Tabellen-Design
%   - etwas mehr Zeilenluft, dezente graue Linien (wirkt auch mit \hline ruhiger)
%   - Kopfzeile einer Tabelle bei Bedarf zart hinterlegen: \rowcolor{bfwTabHead}
% ---------------------------------------------------------------------------
\renewcommand{\arraystretch}{1.25}
\arrayrulecolor{bfwInkSoft!50}
\setlength{\heavyrulewidth}{1pt}
\setlength{\lightrulewidth}{0.5pt}

% ---------------------------------------------------------------------------
% Boxen — klare farbige Titelbalken mit gesperrten Versal-Labels
% (Titel bleibt per Option überschreibbar: \begin{merksatz}[title={...}])
% ---------------------------------------------------------------------------
\tcbset{bfwbox/.style={%
  enhanced, breakable,
  boxrule=0.5pt, arc=2.5pt,
  left=10pt, right=10pt, top=8pt, bottom=8pt,
  toptitle=4pt, bottomtitle=4pt,
  titlerule=0pt,
  fonttitle=\bfwlabelfont,
  coltitle=white
}}

% Merkkästchen — wichtige Definition / Kernpunkt
\newtcolorbox{merksatz}[1][]{%
  bfwbox,
  colback=bfwBgTeal,
  colframe=bfwPrimaryDark,
  colbacktitle=bfwPrimaryDark,
  title={MERKE},
  #1
}

% Eselsbrücke — Mnemoniks
\newtcolorbox{eselsbruecke}[1][]{%
  bfwbox,
  colback=bfwBgRose,
  colframe=bfwAccent!85!black,
  colbacktitle=bfwAccent!85!black,
  title={ESELSBRÜCKE},
  #1
}

% Beispiel-Box
\newtcolorbox{beispiel}[1][]{%
  bfwbox,
  colback=bfwBgSoft,
  colframe=bfwInkSoft,
  colbacktitle=bfwInkSoft,
  title={BEISPIEL},
  #1
}

% Lösungs-Box (für Aufgabenhefte)
\newtcolorbox{loesung}[1][]{%
  bfwbox,
  colback=bfwSuccess!7,
  colframe=bfwSuccess!65!black,
  colbacktitle=bfwSuccess!65!black,
  title={LÖSUNG},
  #1
}

% Achtung-Box — typische Fehler / Stolperfallen
\newtcolorbox{achtung}[1][]{%
  bfwbox,
  colback=bfwWarn!6,
  colframe=bfwWarn!85!black,
  colbacktitle=bfwWarn!85!black,
  title={ACHTUNG},
  #1
}

% Formel-Box — für zentrale Formeln (groß, ruhig, ohne Titelbalken)
\newtcolorbox{formelbox}[1][]{%
  enhanced,
  colback=bfwBgTeal,
  colframe=bfwPrimaryDark,
  boxrule=1pt, arc=3pt,
  left=10pt, right=10pt, top=6pt, bottom=6pt,
  #1
}

% Glossar-Eintrag
\newcommand{\glossareintrag}[2]{%
  \par\noindent\textbf{\color{bfwPrimaryDark}#1} \enspace #2\par\smallskip
}

% ---------------------------------------------------------------------------
% Abschnitts-Design: farbiges Nummern-Quadrat + feine Linie unter \section,
% dezent farbige \subsection/\subsubsection
% ---------------------------------------------------------------------------
\usepackage{titlesec}
\newcommand{\bfwSecBadge}[1]{%
  \begin{tikzpicture}[baseline={([yshift=-0.62em]n.center)}]
    \node[fill=bfwPrimaryDark, text=white, font=\normalsize\bfseries,
          minimum width=1.55em, minimum height=1.55em, inner sep=1pt,
          rounded corners=1.5pt] (n) {#1};
  \end{tikzpicture}}
\titleformat{\section}
  {\Large\bfseries\color{bfwPrimaryDark}}
  {\bfwSecBadge{\thesection}}{0.6em}{}
  [\vspace{-0.2em}{\color{bfwPrimary!60}\titlerule[0.8pt]}]
\titleformat{\subsection}
  {\large\bfseries\color{bfwPrimaryDark}}
  {\textcolor{bfwPrimary}{\thesubsection}}{0.5em}{}
\titleformat{\subsubsection}
  {\normalsize\bfseries\color{bfwInk}}
  {\textcolor{bfwPrimary}{\thesubsubsection}}{0.4em}{}
\titlespacing*{\section}{0pt}{1.6em}{1em}
\titlespacing*{\subsection}{0pt}{1.2em}{0.5em}

% TikZ-Stil "skizzenhaft" — etwas weicher als Rough.js, aber stabil
\tikzset{
  roughbox/.style = {
    draw=bfwPrimaryDark, thick, fill=bfwBgTeal,
    rounded corners=4pt, inner sep=6pt
  },
  roughaccent/.style = {
    draw=bfwAccent, thick, fill=bfwBgAmber,
    rounded corners=4pt, inner sep=6pt
  },
  roughline/.style = {
    draw=bfwInkSoft, thick, -{Stealth[length=2.5mm]}
  }
}

% Bullet-Symbole (bewusst kein FontAwesome — vermeidet Font-Installations-Probleme)
\newcommand{\faLightbulb}{$\bullet$}
\newcommand{\faBookmark}{$\bullet$}

% ---------------------------------------------------------------------------
% Kopf-/Fußzeile
%   Kopf links:  Wortlogo „Wirtschaftsfachwirt"
%   Kopf rechts: Dokument-Kurztext, gesetzt via \kopfzeile{...}
%   Fuß links:   © Can Siebert · 2026 — Fuß rechts: Seite X von Y
%   Titelseite (Seite 1) bleibt ohne Kopf-/Fußzeile (\handouttitel setzt
%   \thispagestyle{empty}).
% ---------------------------------------------------------------------------
\usepackage{fancyhdr}
\newcommand{\bfwKopfzeileText}{}
\newcommand{\kopfzeile}[1]{\renewcommand{\bfwKopfzeileText}{#1}}
\pagestyle{fancy}
\fancyhf{}
\fancyhead[L]{\small\bfseries\color{bfwPrimaryDark}Wirtschaftsfachwirt}
\fancyhead[R]{\small\color{bfwInkSoft}\bfwKopfzeileText}
\renewcommand{\headrulewidth}{0.6pt}
\renewcommand{\headrule}{{\color{bfwPrimary}\hrule width\headwidth height 0.6pt}}
\fancyfoot[L]{\small\color{bfwInkSoft}\copyright\ Can Siebert\,\textperiodcentered\,2026}
\fancyfoot[R]{\small\color{bfwInkSoft}Seite \thepage\ von \pageref*{LastPage}}
% Auch \thispagestyle{plain} (z.B. durch \maketitle) soll leer bleiben:
\fancypagestyle{plain}{\fancyhf{}\renewcommand{\headrulewidth}{0pt}\renewcommand{\headrule}{}}

% ---------------------------------------------------------------------------
% \handouttitel{Obertitel}{Untertitel}{Kurs-Label}{Termin-Zeile}
%   Einheitlicher professioneller Titelblock: Dachzeile, farbige Headline,
%   Untertitel, farbige Trennlinie und dezente Meta-Zeile (Kurs/Termin/Dozent).
%   Ersetzt \title/\maketitle. Direkt nach \begin{document} aufrufen.
% ---------------------------------------------------------------------------
\newcommand{\handouttitel}[4]{%
  \thispagestyle{empty}%
  \begingroup
  \setlength{\parindent}{0pt}%
  {\bfwlabelfont\color{bfwPrimary}WIRTSCHAFTSFACHWIRT (IHK)\par}
  \vspace{0.55em}
  {\huge\bfseries\color{bfwPrimaryDark}#1\par}
  \vspace{0.5em}
  {\large\color{bfwInkSoft}#2\par}
  \vspace{0.9em}
  {\color{bfwPrimary}\rule{\linewidth}{2pt}}\par
  \vspace{0.55em}
  {\small\color{bfwInkSoft}#3\ \,\textperiodcentered\,\ #4\ \,\textperiodcentered\,\ Dozent: Can Siebert\par}
  \endgroup
  \vspace{1.4em}%
}


\kopfzeile{Handout BWL Lektion 2 \textperiodcentered{} Teil 2}

\begin{document}
\handouttitel{Lektion~2 \textperiodcentered{} Teil~2: Controlling, Personal \& Zusammenwirken}%
  {Controlling, Personalwirtschaft --- und wie alle betrieblichen Funktionen zusammenwirken}%
  {1.2 Betriebliche Funktionen}%
  {Sa, 12.09.2026 \textperiodcentered{} Session 2}

\begin{merksatz}[title={\bfseries Worum geht es in dieser Session?}]
In Teil~1 (heute Vormittag) haben Sie das \emph{Rechnungswesen}, \emph{Investition und
Finanzierung} sowie die wichtigsten \emph{Kennzahlen-Formeln} kennengelernt. Teil~2
schließt die Lektion ab: Das \emph{Controlling} plant und steuert die
Unternehmensprozesse --- mit den Kennzahlen aus Teil~1 als wichtigstem Werkzeug ---, die
\emph{Personalwirtschaft} stellt die Mitarbeiterinnen und Mitarbeiter bereit. Zum Schluss
betrachten wir das \emph{Zusammenwirken} aller betrieblichen Funktionen (Rahmenplan
1.2.2): Nur wenn alle Bereiche verzahnt arbeiten --- und ihre Zielkonflikte gelöst
werden ---, erreicht das Unternehmen seine Ziele.
\end{merksatz}

\tableofcontents
\vspace{1em}
\hrule
\vspace{1em}

\section{Controlling}

\subsection{Begriff: mehr als Kontrolle}

\textbf{Controlling} (von engl. \emph{to control} = steuern, lenken) ist eine
\textbf{Management-Service-Funktion} zur Planung, Steuerung und Koordination der
Unternehmensprozesse. Es geht weit über den deutschen Begriff „Kontrolle`` hinaus:
\emph{Kontrolle} ist der vergangenheitsbezogene Soll-Ist-Vergleich --- \emph{Controlling}
umfasst zusätzlich Planung, Informationsversorgung und die zukunftsgerichtete Steuerung.
Der Controller ist Berater und „Lotse`` des Managements, nicht dessen Kontrolleur.
Prüfungsklassiker: Abgrenzung zur \emph{internen Revision} --- sie prüft
vergangenheitsorientiert und unabhängig die Ordnungsmäßigkeit, während das Controlling
zukunftsgerichtet plant und steuert.

\subsection{Der Controlling-Kreislauf}

\begin{itemize}
  \item \textbf{Planung:} Ziele setzen, Maßnahmen und Budgets festlegen (Sollgrößen).
  \item \textbf{Kontrolle:} Soll-Ist-Vergleich, Abweichungsanalyse (Ursachen ermitteln).
  \item \textbf{Steuerung:} Gegensteuerungsmaßnahmen einleiten, Pläne anpassen.
  \item \textbf{Information:} entscheidungsrelevante Informationen beschaffen, aufbereiten
    und den Führungskräften bereitstellen (Berichtswesen/Reporting).
\end{itemize}

\subsection{Strategisches vs. operatives Controlling}

\begin{center}
\begin{tabular}{p{3.2cm} p{5.2cm} p{5.2cm}}
 & \textbf{Strategisches Controlling} & \textbf{Operatives Controlling}\\
\hline
Leitfrage & „Tun wir die richtigen Dinge?`` & „Tun wir die Dinge richtig?``\\
Zeithorizont & langfristig ($>$ 5 Jahre) & kurz- bis mittelfristig (bis ca. 3 Jahre)\\
Orientierung & Umwelt und Unternehmen: Chancen/Risiken, Stärken/Schwächen & innerbetrieblich:
Aufwand/Ertrag, Kosten/Leistungen\\
Zielgrößen & Erfolgspotenziale, Existenzsicherung & Gewinn, Rentabilität, Liquidität,
Wirtschaftlichkeit\\
Instrumente & SWOT-Analyse, Portfolio, Szenariotechnik & Budgetierung, Soll-Ist-Vergleich,
Kennzahlen, Deckungsbeitragsrechnung\\
\end{tabular}
\end{center}

\subsection{Kennzahlen als Controlling-Werkzeug}

Kennzahlen verdichten betriebliche Sachverhalte zu aussagekräftigen Zahlen. Wichtige
Beispiele: \textbf{Eigenkapitalquote} (Kapitalstruktur), \textbf{Eigenkapital-,
Gesamtkapital- und Umsatzrentabilität} (Erfolg), \textbf{Liquiditätsgrade}
(Zahlungsfähigkeit) und \textbf{Cashflow} (Innenfinanzierungskraft). Kennzahlen
entfalten ihren Wert erst im \emph{Vergleich}: Zeitvergleich, Betriebsvergleich,
Soll-Ist-Vergleich. Alle Formeln mit Rechenweg und einfachen Rechenbeispielen finden
Sie im \textbf{Handout zu Teil~1} (Abschnitt „Formeln \& Rechenbeispiele``).

\begin{beispiel}
Das operative Controlling der Nordwest Küchengeräte AG meldet: Die Umsatzrentabilität
ist im Zeitvergleich von 7~\% auf 5~\% gesunken. Der Soll-Ist-Vergleich zeigt
Kostenüberschreitungen in der Fertigung --- das Controlling schlägt als Steuerungsmaßnahme
eine Nachkalkulation der Produktlinien und ein straffes Kostenbudget vor.
\end{beispiel}

\section{Personalwirtschaft}

Im globalen Wettbewerb gilt der Inputfaktor Arbeit als entscheidender strategischer
Erfolgsfaktor: Rekrutierung, Bindung und Motivation der Mitarbeiter zählen zu den
wichtigsten Aufgaben des Managements.

\subsection{Aufgabenbereiche der Personalwirtschaft}

\begin{itemize}
  \item \textbf{Personalplanung:} den künftigen Personalbedarf nach Menge, Qualifikation,
    Zeit und Ort ermitteln.
  \item \textbf{Personalbeschaffung:} Stellen intern (Versetzung, Aufstieg) oder extern
    (Anzeigen, Jobportale, Personalberater, Arbeitsagentur) besetzen; Auswahl z.\,B. über
    Interviews und Assessmentcenter.
  \item \textbf{Personalentwicklung:} Aus- und Weiterbildung, Konzeption von
    Weiterbildungsangeboten, Mitarbeitergespräche, Potenzialanalysen
    (z.\,B. Mitarbeiter-Portfolio: Stars, Arbeitstiere, Problemfälle, Leistungsschwache).
  \item \textbf{Personaleinsatz:} Zuordnung der Mitarbeiter zu Arbeitsplätzen,
    Arbeitsplatzbeschreibungen, ergonomische Arbeitsplatzgestaltung.
  \item \textbf{Personalfreisetzung:} Abbau personeller Überkapazitäten --- möglichst
    sozialverträglich (Fluktuation nutzen, Altersteilzeit, Abfindungen, erst zuletzt
    Kündigungen).
  \item Ergänzend: Personal\emph{verwaltung} (Personalakten), -\emph{betreuung}
    (Sozialleistungen), -\emph{entlohnung}, -\emph{beurteilung} (Zeugnisse) und
    Personal\emph{controlling} (z.\,B. Personalkosten- und Fehlzeitenanalyse).
\end{itemize}

\begin{eselsbruecke}
Der Personal-Lebenszyklus: \textbf{P}lanen $\rightarrow$ \textbf{B}eschaffen $\rightarrow$
\textbf{E}ntwickeln $\rightarrow$ \textbf{E}insetzen $\rightarrow$ \textbf{F}reisetzen ---
„\textbf{P}ersonal \textbf{b}raucht \textbf{e}ine \textbf{e}chte \textbf{F}ührung.``
\end{eselsbruecke}

\subsection{Ziele der Personalwirtschaft}

\begin{itemize}
  \item \textbf{Wirtschaftliche Ziele:} die richtigen Mitarbeiter in der richtigen Anzahl
    mit der richtigen Qualifikation zur richtigen Zeit am richtigen Ort --- bei optimalen
    Personalkosten und hoher Arbeitsproduktivität.
  \item \textbf{Soziale Ziele:} Zufriedenheit, Motivation und Gesundheit der Mitarbeiter,
    gerechte Entlohnung, Entwicklungsmöglichkeiten, Vereinbarkeit von Familie und Beruf.
\end{itemize}

\begin{beispiel}
\textbf{Demografischer Wandel:} Die Bevölkerung altert, die Erwerbsbevölkerung schrumpft.
Folgen für Betriebe: Fachkräftemangel, alternde Belegschaften, Verlust von Erfahrungswissen.
Gegenmaßnahmen: Ausbildung und Übernahme, lebenslanges Lernen, Gesundheitsmanagement,
altersgemischte Teams, familienfreundliche Angebote (z.\,B. Betriebskita), Rekrutierung
neuer Zielgruppen (Ältere, Zuwanderer, Frauen in Teilzeit), Employer Branding.
\end{beispiel}

\section{Zusammenwirken der betrieblichen Funktionen (1.2.2)}

\subsection{Gemeinsames Ziel, verzahnte Bereiche}

Alle betrieblichen Teilbereiche müssen ihre Aufgaben erfüllen, damit das Unternehmen als
Ganzes rentabel arbeitet und seine Ziele erreicht. Die Funktionsbereiche sind über
\textbf{Schnittstellen} miteinander verzahnt: Der \emph{Realgüterstrom} fließt von den
Beschaffungsmärkten durch das Unternehmen zu den Absatzmärkten, der
\emph{Nominalgüterstrom} (Geld) in umgekehrter Richtung --- gesteuert durch eine
zielorientierte Unternehmensführung mit Controlling und verbunden über
Informationssysteme mit Staat, Geld- und Kapitalmarkt und natürlicher Umwelt.

\subsection{Die Wertschöpfungskette}

Nach Porter lassen sich die Aktivitäten eines Unternehmens als
\textbf{Wertschöpfungskette} darstellen:
\begin{itemize}
  \item \textbf{Primäre Aktivitäten} (unmittelbar wertschöpfend): Eingangslogistik
    $\rightarrow$ Produktion $\rightarrow$ Ausgangslogistik $\rightarrow$
    Marketing/Vertrieb $\rightarrow$ Kundendienst.
  \item \textbf{Unterstützende Aktivitäten:} Unternehmensinfrastruktur (u.\,a.
    Rechnungswesen, Finanzierung, Controlling), Personalwirtschaft, Technologieentwicklung,
    Beschaffung.
\end{itemize}
Wertschöpfung = Wert der erstellten Leistung abzüglich der von außen bezogenen
Vorleistungen. Jede Funktion soll zur Wertschöpfung beitragen --- Schnittstellen sind
dabei potenzielle Reibungsverluste.

\subsection{Zielkonflikte zwischen Funktionsbereichen}

\begin{center}
\begin{tabular}{p{4.4cm} p{4.4cm} p{5.0cm}}
\textbf{Bereich A will …} & \textbf{Bereich B will …} & \textbf{Konflikt}\\
\hline
Absatz/Marketing: große Produktvielfalt, kurze Lieferzeiten, hohe Lieferbereitschaft &
Produktion: Standardisierung, große Lose, gleichmäßige Auslastung & Variantenvielfalt
erhöht Rüstzeiten und Komplexität\\
Einkauf: große Bestellmengen (Mengenrabatte) & Finanzen/Lager: geringe Kapitalbindung,
niedrige Lagerkosten & Rabattvorteil vs. Lager- und Zinskosten\\
Vertrieb: großzügige Zahlungsziele für Kunden & Finanzen: schneller Zahlungseingang,
Liquidität & Umsatzförderung vs. Forderungsrisiko\\
Personal: Weiterbildung, familienfreundliche Leistungen & Controlling: Kostensenkung,
Budgeteinhaltung & Investition in Mitarbeiter vs. kurzfristige Kosten\\
\end{tabular}
\end{center}

\begin{merksatz}
Zielkonflikte sind normal --- entscheidend ist ihre \emph{Koordination}: durch die
Unternehmensführung, das Controlling (Planung, Budgets, Kennzahlen) und abgestimmte
Ziele entlang der Wertschöpfungskette. Prüfungsklassiker: „Absatz will Vielfalt,
Produktion will Standardisierung.``
\end{merksatz}

\begin{merksatz}[title={\bfseries Wichtig für die Prüfungsvorbereitung (Teil 2)}]
Sie sollten sicher können: die Abgrenzung \emph{Controlling/Kontrolle} (und zur internen
Revision) samt \emph{Controlling-Kreislauf} (Planung--Kontrolle--Steuerung--Information),
die Unterscheidung \emph{strategisches vs. operatives Controlling}, den Einsatz von
\emph{Kennzahlen} im Vergleich (Formeln: siehe Handout Teil~1), die
\emph{Aufgabenbereiche und Ziele der Personalwirtschaft} (Planung, Beschaffung,
Entwicklung, Einsatz, Freisetzung) samt demografischem Wandel sowie typische
\emph{Zielkonflikte} zwischen Funktionsbereichen und die \emph{Wertschöpfungskette}.
Üben Sie mit dem Aufgabenheft und dem Quiz zu Teil~2!
\end{merksatz}

\vspace{1em}
\hrule
\vspace{0.8em}
\noindent\textsf{\small Quellen: Rahmenplan Wirtschaftsbezogene Qualifikationen (DIHK),
Abschnitte 1.2.1.6--1.2.1.7 und 1.2.2; Übungsband Betriebswirtschaft (DIHK-Bildungs-GmbH),
Kap.~1, Übungen 25--44. Der Textband-Auszug wird nachgereicht. Nächste Lektion
(Sa, 26.09.2026): Existenzgründung \& Rechtsformen.}

\end{document}
